% ===========================================================================
% MANUAL DE USUARIO - TEST VOCACIONAL
% Proyecto: test-vocacional
% Destinatario: Usuarios finales del sistema
% Fecha: Abril 2025
% ===========================================================================

\documentclass[12pt, a4paper]{report}

% ---------------------------------------------------------------------------
% PAQUETES BÁSICOS
% ---------------------------------------------------------------------------
\usepackage[utf8]{inputenc}
\usepackage[T1]{fontenc}
\usepackage[spanish]{babel}
\usepackage{geometry}
\usepackage{graphicx}
\usepackage{hyperref}
\usepackage{listings}
\usepackage{xcolor}
\usepackage{fancyhdr}
\usepackage{titlesec}
\usepackage{tocloft}
\usepackage{tabularx}
\usepackage{longtable}
\usepackage{multirow}
\usepackage{booktabs}
\usepackage{amsmath}
\usepackage{amssymb}
\usepackage{enumitem}
\usepackage{verbatim}
\usepackage{float}
\usepackage{tikz}
\usepackage{pgfplots}
\usepackage{minted}
\usepackage{caption}
\usepackage{subcaption}
\usepackage{pifont}        % Para símbolos checkmark/X
\usepackage{awesomebox}    % Para iconos (nota, advertencia, peligro)
\usepackage{tcolorbox}
\usepackage{wrapfig}

% ---------------------------------------------------------------------------
% CONFIGURACIÓN DE PÁGINA
% ---------------------------------------------------------------------------
\geometry{
    top=3cm,
    bottom=2.5cm,
    left=2.5cm,
    right=2.5cm
}

\pagestyle{fancy}
\fancyhf{}
\fancyhead[L]{\leftmark}
\fancyhead[R]{\thepage}
\fancyfoot[C]{}
\renewcommand{\headrulewidth}{0.5pt}

% ---------------------------------------------------------------------------
% COLORES PERSONALIZADOS
% ---------------------------------------------------------------------------
\definecolor{primaryblue}{RGB}{0, 82, 147}
\definecolor{successgreen}{RGB}{34, 139, 34}
\definecolor{warningorange}{RGB}{255, 140, 0}
\definecolor{dangerred}{RGB}{220, 20, 60}
\definecolor{infoblue}{RGB}{30, 144, 255}
\definecolor{lightgray}{RGB}{240, 240, 240}

% ---------------------------------------------------------------------------
% LISTINGS CONFIGURACIÓN
% ---------------------------------------------------------------------------
\lstdefinestyle{codeStyle}{
    backgroundcolor=\color{lightgray},
    basicstyle=\ttfamily\footnotesize,
    keywordstyle=\color{primaryblue},
    commentstyle=\color{successgreen},
    stringstyle=\color{dangerred},
    frame=single,
    framerule=0.5pt,
    breaklines=true,
    numbers=left,
    numberstyle=\tiny\color{gray}
}

\lstset{style=codeStyle}

% ---------------------------------------------------------------------------
% TCOLORBOX - CAJAS DE NOTAS
% ---------------------------------------------------------------------------
\tcbset{colback=white, colframe=primaryblue, fonttitle=\bfseries}

% ---------------------------------------------------------------------------
% META DATOS PDF
% ---------------------------------------------------------------------------
\hypersetup{
    colorlinks=true,
    linkcolor=primaryblue,
    filecolor=magenta,
    urlcolor=infoblue,
    pdftitle={Manual de Usuario - Test Vocacional},
    pdfauthor={Equipo de Desarrollo},
    pdfsubject={Manual de Usuario},
    pdfkeywords={Test Vocacional, Manual de Usuario, RIASEC}
}

% ===========================================================================
% COMANDOS PERSONALIZADOS
% ===========================================================================
\newcommand{\roleadmin}{\textcolor{dangerred}{\textbf{Administrador}}}
\newcommand{\rolestudent}{\textcolor{successgreen}{\textbf{Estudiante}}}
\newcommand{\roledece}{\textcolor{primaryblue}{\textbf{DECE}}}
\newcommand{\rolezonal}{\textcolor{warningorange}{\textbf{Zonal}}}
\newcommand{\required}{\textcolor{dangerred}{*}}
\newcommand{\stepicon}{\ding{51}}
\newcommand{\warningicon}{\ding{72}}
\newcommand{\important}{\tcblower\textbf{Nota:}}

% ===========================================================================
% INICIO DEL DOCUMENTO
% ===========================================================================

\begin{document}

% ---------------------------------------------------------------------------
% PORTADA
% ---------------------------------------------------------------------------
\begin{titlepage}
    \centering
    \vspace*{2cm}
    
    {\Huge \textbf{MANUAL DE USUARIO}}\\[0.5cm]
    {\Huge \textbf{TEST VOCACIONAL}}\\[2cm]
    
    \includegraphics[width=0.4\textwidth]{./logo.png}\\[1cm]
    
    {\Large \textbf{Sistema de Orientación Vocacional}}\\[0.5cm]
    {\large Modelo RIASEC - Holland}\\[1.5cm]
    
    \rule{\textwidth}{0.4pt}\\[1cm]
    {\large \textbf{Guía completa para usuarios}}\\
    {\large Estudiantes · DECE · Zonal · Administrador}\\[1cm]
    \rule{\textwidth}{0.4pt}\\[2cm]
    
    {\large Versión 1.0}\\[0.5cm]
    {\large Abril 2025}\\[3cm]
    
    {\normalsize \textbf{Contenido:}}\\
    {\normalsize Guías paso a paso, capturas de pantalla,}\\
    {\normalsise preguntas frecuentes y procedimientos.}
    
    \vfill
\end{titlepage}

% ---------------------------------------------------------------------------
% PÁGINA DE AUTORES/CONTACTO
% ---------------------------------------------------------------------------
\newpage
\thispagestyle{empty}
\vspace*{5cm}
\begin{flushleft}
    {\large \textbf{Información de Contacto}}\\[0.5cm]
    \textbf{Soporte Técnico:}\\
    Email: soporte@testvocacional.edu.ec\\
    Teléfono: +593 2 XXX-XXXX\\[1cm]
    
    \textbf{Desarrollo del Sistema:}\\
    Equipo de Desarrollo - test-vocacional\\
    GitHub: https://github.com/[org]/test-vocacional\\[2cm]
    
    {\large \textbf{Ubicación del Proyecto}}\\[0.5cm]
    \texttt{/Users/rubenjaramillo/test-vocacional2} \\
\end{flushleft}

% ---------------------------------------------------------------------------
% RESUMEN
% ---------------------------------------------------------------------------
\chapter*{RESUMEN}
\addcontentsline{toc}{chapter}{Resumen}

Este Manual de Usuario describe el uso del \textbf{Sistema de Test Vocacional}, una herramienta en línea basada en el modelo RIASEC de Holland para orientar a estudiantes en su elección de carrera.

\textbf{¿A quién está dirigido?}
\begin{itemize}
    \item \textbf{Estudiantes}: Realizan el test y consultan sus resultados
    \item \textbf{Personal DECE}: Gestionan y monitorean estudiantes de su institución
    \item \textbf{Coordinadores Zonales}: Supervisan resultados por zonas educativas
    \item \textbf{Administradores}: Configuran el sistema y gestionan contenido
\end{itemize}

\textbf{Contenido del manual:}
\begin{enumerate}
    \item Introducción al sistema y modelo RIASEC
    \item Acceso y configuración inicial
    \item Guías por rol (Estudiante, DECE, Zonal, Admin)
    \item Procesos paso a paso con capturas ilustrativas
    \item Preguntas frecuentes y solución de problemas
    \item Glosario de términos técnicos
\end{enumerate}

\section*{Cómo usar este manual}
Este manual está organizado por \textbf{roles de usuario}. Comience por el capítulo correspondiente a su función en el sistema. Cada sección incluye:
\begin{itemize}
    \item \textbf{Objetivo}: Qué logrará con esa funcionalidad
    \item \textbf{Paso a paso}: Instrucciones numbered list
    \item \textbf{Importante}: Notas y consideraciones
    \item \textbf{Visual}: Descripción de lo que verá en pantalla
\end{itemize}

% ---------------------------------------------------------------------------
% ÍNDICE
% ---------------------------------------------------------------------------
\tableofcontents
\newpage

\listoffigures
\newpage

% ===========================================================================
% CAPÍTULO 1: INTRODUCCIÓN
% ===========================================================================
\chapter{INTRODUCCIÓN AL SISTEMA}

\section{¿Qué es el Test Vocacional?}
El Test Vocacional es una herramienta de orientación profesional que ayuda a estudiantes a identificar sus intereses, habilidades y valores laborales mediante el \textbf{modelo RIASEC} desarrollado por el psicólogo John L. Holland.

\section{El Modelo RIASEC}
RIASEC representa seis categorías vocacionales:

\begin{table}[htbp]
\centering
\begin{tabularx}{\textwidth}{|c|X|l|}
\hline
\textbf{Código} & \textbf{Tipo} & \textbf{Descripción} \\
\hline
\textcolor{red}{R} & \textbf{Realista} & Personas prácticas, que disfrutan trabajos manuales y físicos. Ej: ingeniería, construcción, mecánica. \\
\hline
\textcolor{blue}{I} & \textbf{Investigador} & Personas analíticas, científicas, curiosas. Ej: medicina, investigación, tecnología. \\
\hline
\textcolor{purple}{A} & \textbf{Artístico} & Personas creativas, expresivas, innovadoras. Ej: diseño, música, arquitectura. \\
\hline
\textcolor{green}{S} & \textbf{Social} & Personas que ayudan, enseñan, cuidan. Ej: docencia, psicología, enfermería. \\
\hline
\textcolor{orange}{E} & \textbf{Emprendedor} & Personas persuasivas, líderes, empresarias. Ej: administración, derecho, ventas. \\
\hline
\textcolor{cyan}{C} & \textbf{Convencional} & Personas organizadas, detallistas, administrativas. Ej: contabilidad, secretariado, logística. \\
\hline
\end{tabularx}
\caption{Categorías del modelo RIASEC}
\end{table}

\section{¿Cómo funciona?}
\begin{enumerate}
    \item \textbf{Encuesta previa}: Recopila información sobre preferencias, familia y actividades
    \item \textbf{Cuestionario}: 60 preguntas (10 por categoría) con escala Likert 1-5
    \item \textbf{Cálculo automático}: El sistema suma puntajes por categoría
    \item \textbf{Resultados}: Muestra gráficos, recomendaciones y carreras sugeridas
\end{enumerate}

\section{Acceso al Sistema}
\begin{itemize}
    \item \textbf{URL}: \texttt{http://localhost/test-vocacional} (desarrollo) o dominio institucional
    \item \textbf{Navegadores compatibles}: Chrome, Firefox, Safari, Edge (última versión)
    \item \textbf{Dispositivos}: Computadora, tablet (se recomienda pantalla grande para mejor experiencia)
\end{itemize}

% ===========================================================================
% CAPÍTULO 2: ACCESO Y CUENTA
% ===========================================================================
\chapter{ACCESO Y GESTIÓN DE CUENTA}

\section{Primer Acceso - Registro}

\subsection{¿Ya tienes cuenta?}
Si ya te han registrado (el administrador te creó una cuenta):
\begin{enumerate}
    \item Ve a la página de \textbf{Login}
    \item Ingresa tu \textbf{usuario} (puede ser tu cédula o email)
    \item Ingresa tu \textbf{contraseña}
    \item Marca \textbf{"Recordarme"} si es un dispositivo personal
    \item Haz clic en \textbf{"Iniciar Sesión"}
\end{enumerate}

\subsection{¿Primera vez? - Registro}
Si no tienes cuenta, primero debes ser registrado por un administrador o DECE de tu institución. \textbf{No puedes auto-registrarte}. Contacta a:
\begin{itemize}
    \item Tu profesor/asesor vocacional
    \item El departamento DECE de tu institución
    \item El administrador del sistema
\end{itemize}

Ellos te proporcionarán:
\begin{itemize}
    \itemUsuario (generalmente tu número de cédula)
    \item Contraseña temporal (que deberás cambiar al primer login)
\end{itemize}

\section{Recuperación de Contraseña}
Si olvidaste tu contraseña:
\begin{enumerate}
    \item En la página de login, haz clic en \textbf{¿Olvidaste tu contraseña?}
    \item Ingresa el \textbf{email} registrado en tu cuenta
    \item Revisa tu correo (bandeja de entrada y spam)
    \item Haz clic en el \textbf{enlace de recuperación} (válido por 24 horas)
    \item Crea una nueva contraseña (mínimo 6 caracteres)
    \item Inicia sesión con la nueva contraseña
\end{enumerate}

\section{Cambiar Contraseña}
Una vez logueado:
\begin{enumerate}
    \item Haz clic en tu \textbf{nombre} (esquina superior derecha)
    \item Selecciona \textbf{"Mi Perfil"}
    \item Busca la sección \textbf{"Cambiar Contraseña"}
    \item Ingresa tu contraseña actual
    \item Ingresa la nueva contraseña (2 veces)
    \item Haz clic en \textbf{"Guardar"}
    \item Verás un mensaje de confirmación verde
\end{enumerate}

\section{Requisitos de Contraseña}
\begin{table}[htbp]
\centering
\begin{tabularx}{\textwidth}{|l|X|}
\hline
\textbf{Requisito} & \textbf{Detalle} \\
\hline
Longitud mínima & 6 caracteres \\
Caracteres permitidos & Letras, números y símbolos \\
Cambio forzado & Debes cambiar la contraseña temporal al primer acceso \\
\hline
\end{tabularx}
\end{table}

\section{Cerrar Sesión}
\begin{itemize}
    \item Haz clic en \textbf{"Cerrar Sesión"} en el menú superior
    \item O simplemente cierra el navegador (la sesión expira en 1 hora)
\end{itemize}

\section{Perfil de Usuario}
En la página de \textbf{Perfil} puedes:
\begin{itemize}
    \item Ver tu información personal
    \item Actualizar tu foto (si está habilitado)
    \item Cambiar tu contraseña
    \item Ver tu historial de tests
\end{itemize}

% ===========================================================================
% CAPÍTULO 3: GUÍA PARA ESTUDIANTES
% ===========================================================================
\chapter{GUÍA PARA ESTUDIANTES (\rolestudent)}

\section{¿Qué puedes hacer como Estudiante?}
Como estudiante, tienes acceso a:
\begin{itemize}
    \item Realizar el test vocacional (una vez cada 6 meses)
    \item Consultar tus resultados anteriores
    \item Descargar tu reporte individual en PDF
    \item Actualizar tu información básica
\end{itemize}

\section{Proceso 1: Completar el Test Vocacional}

\subsection{Paso 1: Acceso al test}
\begin{enumerate}
    \item Inicia sesión con tu usuario y contraseña
    \item En el \textbf{Dashboard} o página principal, verás un botón grande: \textbf{"Comenzar Test Vocacional"}
    \item Haz clic en ese botón
\end{enumerate}

\subsection{Paso 2: Encuesta previa (OBLIGATORIA)}
Antes del test, debes completar una breve encuesta. \textbf{No puedes omitirla}.

\textbf{Secciones de la encuesta:}
\begin{enumerate}
    \item \textbf{Profesiones preferidas}: Selecciona 3 profesiones que más te gustaría ejercer
    \item \textbf{Profesiones menos preferidas}: Selecciona 3 que NO te gustaría
    \item \textbf{Datos de la madre}: Nivel de estudios y profesión
    \item \textbf{Datos del padre}: Nivel de estudios y profesión
    \item \textbf{Actividades de tiempo libre}: 3 actividades que disfrutas
    \item \textbf{Características de éxito}: 3 cualidades que consideras importantes
    \item \textbf{Importancia del éxito}: Describe qué significa para ti
\end{enumerate}

\textbf{Importante}: Todos los campos son obligatorios. Si omites uno, el sistema no te permitirá continuar.

\subsection{Paso 3: Cuestionario RIASEC (60 preguntas)}
Ahora verás el cuestionario principal. \textbf{Características:}
\begin{itemize}
    \item \textbf{Formato}: Preguntas presentadas como un "libro" con páginas
    \item \textbf{Barra de progreso}: Muestra cuánto has avanzado (arriba)
    \item \textbf{Navegación}: Botones \textbf{« Anterior} y \textbf{Siguiente »}
    \item \textbf{Respuestas}: Escala Likert del 1 al 5
    \begin{itemize}
        \item \textbf{1}: Totalmente en desacuerdo
        \item \textbf{2}: En desacuerdo
        \item \textbf{3}: Neutral / Ni de acuerdo ni en desacuerdo
        \item \textbf{4}: De acuerdo
        \item \textbf{5}: Totalmente de acuerdo
    \end{itemize}
\end{itemize}

\textbf{Flujo de respuestas:}
\begin{enumerate}
    \item Lee cada afirmación cuidadosamente
    \item Selecciona el número que mejor refleje tu opinión
    \item Haz clic en \textbf{"Siguiente"} para continuar
    \item Puedes volver atrás (\textbf{« Anterior}) para cambiar respuestas
\end{enumerate}

\textbf{Guardado automático}: Tus respuestas se guardan en el navegador cada 30 segundos. Si cierras la ventana accidentalmente, al volver verás un botón \textbf{"Recuperar respuestas guardadas"}.

\subsection{Paso 4: Finalizar y ver resultados}
\begin{enumerate}
    \item Al responder las 60 preguntas, el botón cambia a \textbf{"Finalizar y Ver Resultados"}
    \item Haz clic para enviar tus respuestas
    \item El sistema calculará tus puntajes (toma unos segundos)
    \item Serás redirigido automáticamente a la página de \textbf{resultados}
\end{enumerate}

\textbf{Importante}: Una vez finalizado, \underline{no puedes volver a realizar el test durante 6 meses}.

\section{Proceso 2: Consultar Resultados}

\subsection{Vista principal de resultados}
Al finalizar el test (o al entrar a \textbf{/results}), verás:

\begin{figure}[H]
\centering
\begin{tcolorbox}[colback=lightgray, colframe=primaryblue, title=Vista de Resultados - Elementos Clave]
\begin{itemize}
    \item \textbf{Gráfico de radar}: Muestra tu perfil en las 6 áreas RIASEC
    \item \textbf{Gráfico de barras}: Puntajes específicos (0-100\%)
    \item \textbf{Tarjetas de categoría}: Cada área con su estado
    \begin{itemize}
        \item \textcolor{successgreen}{APTO (verde)}: ≥80\% - Excelente aptitud
        \item \textcolor{warningorange}{POTENCIAL (amarillo)}: 60-79\% - Buen potencial
        \item \textcolor{dangerred}{POR REFORZAR (rojo)}: <60\% - En desarrollo
    \end{itemize}
    \item \textbf{Tus 3 áreas principales}: Listado con explicación
    \item \textbf{Carreras sugeridas}: Para universidades públicas y privadas
    \item \textbf{Recomendaciones personalizadas} basadas en tu perfil
    \item \textbf{Botón "Descargar PDF"} (Reporte individual)
\end{itemize}
\end{tcolorbox}
\end{figure}

\subsection{Interpretar tus resultados}
\textbf{Estado APTO (verde):}
\begin{itemize}
    \item Tienes una \underline{alta compatibilidad} con esa área
    \item Considera carreras de esa categoría como primera opción
    \item Es una \textbf{fortaleza} destacada
\end{itemize}

\textbf{Estado POTENCIAL (amarillo):}
\begin{itemize}
    \item Tienes \underline{buen potencial} pero puedes desarrollarlo más
    \item Puedes considerar carreras de esa área como segunda opción
    \item Requiere \textbf{refuerzo} en algunas habilidades
\end{itemize}

\textbf{Estado POR REFORZAR (rojo):}
\begin{itemize}
    \item Es un área de \underline{menor afinidad}
    \item Puedes desarrollarla con práctica y formación
    \item No es recomendable como primera opción de carrera
\end{itemize}

\subsection{Descargar tu reporte}
\begin{enumerate}
    \item En la página de resultados, busca el botón azul: \textbf{"Descargar Reporte PDF"}
    \item Haz clic en el botón
    \item El archivo se descargará automáticamente
    \item El PDF incluye:
    \begin{itemize}
        \item Tus datos personales
        \item Gráficos de resultados
        \item Lista de carreras recomendadas
        \item \textbf{Código QR} para validación institucional
        \item Fecha y hora del test
    \end{itemize}
\end{enumerate}

\textbf{Importante}: Guarda este PDF. Puedes enviarlo a tu escuela, universidad o asesor vocacional.

\section{Preguntas Frecuentes - Estudiante}

\textbf{Q: ¿Puedo realizar el test más de una vez?}\\
A: No. El sistema permite \underline{un test cada 6 meses}. Esto asegura que los resultados reflejen un momento de reflexión genuino. Si intentas antes, verás un mensaje indicando cuántos días faltan.

\textbf{Q: ¿Puedo cambiar mis respuestas después de finalizar?}\\
A: No. Una vez enviado el test, las respuestas se guardan y no se pueden modificar. Asegúrate de revisar antes de hacer clic en "Finalizar".

\textbf{Q: ¿Mis resultados son privados?}\\
A: Sí. Solo tú, tu DECE y el personal autorizado pueden verlos. No se comparten fuera de tu institución.

\textbf{Q: ¿Qué pasa si cierro el navegador a mitad del test?}\\
A: Puedes recuperar tus respuestas guardadas. Al volver al test, busca el botón \textbf{"Recuperar respuestas guardadas"}.

\section{Consejos para el Test}
\begin{tcolorbox}[colback=infoblue!10, colframe=infoblue, title=Consejos Importantes]
\begin{itemize}
    \item \textbf{No apresurarse}: Lee cada afirmación con calma
    \item \textbf{Sé honesto}: No hay respuestas "buenas" o "malas"
    \item \textbf{Piensa en ti mismo}: No en lo que otros esperan
    \item \textbf{No dejes preguntas en blanco}: Todas son obligatorias
    \item \textbf{Contexto}: Responde considerando tu situación \underline{actual}
\end{itemize}
\end{tcolorbox}

% ===========================================================================
% CAPÍTULO 4: GUÍA PARA DECE
% ===========================================================================
\chapter{GUÍA PARA DECE (\roledece)}

\section{¿Qué es un DECE?}
\textbf{DECE} = Departamento de Estudios y Consejería Estudiantil. Es el personal encargado del bienestar y orientación de los estudiantes en una institución educativa.

\section{¿Qué puedes hacer como DECE?}
\begin{itemize}
    \item Dashboards institucionales con estadísticas
    \item Ver y gestionar estudiantes de \underline{tu institución únicamente}
    \item Generar reportes grupales por curso o paralelo
    \item Exportar datos a Excel
    \item Descargar reportes individuales de estudiantes
    \item Asignar roles a usuarios de tu institución
    \item Seguimiento vocacional de estudiantes
\end{itemize}

\textbf{Limitación}: Solo ves estudiantes de \underline{tu institución asignada}. No puedes acceder a datos de otras instituciones.

\section{Acceso al Panel DECE}
\begin{enumerate}
    \item Inicia sesión con tus credenciales DECE
    \item Serás redirigido automáticamente al \textbf{Dashboard DECE}
    \item O ve a \texttt{/admin/dece} manualmente
\end{enumerate}

\section{Dashboard DECE - Vista General}
El dashboard muestra:
\begin{table}[htbp]
\centering
\begin{tabularx}{\textwidth}{|l|X|}
\hline
\textbf{Elemento} & \textbf{Descripción} \\
\hline
Estadísticas rápidas & Total estudiantes, tests realizados, promedio general \\
\hline
Gráfico de distribución & Cantidad de estudiantes por categoría RIASEC \\
\hline
Top carreras sugeridas & Las 5 carreras más recomendadas en tu institución \\
\hline
Últimos tests & Lista de estudiantes que realizaron test recientemente \\
\hline
Accesos rápidos & Botones para reportes y gestión \\
\hline
\end{tabularx}
\end{table}

\section{Proceso 1: Gestionar Estudiantes}

\subsection{Ver lista de estudiantes}
\begin{enumerate}
    \item En el menú lateral, haz clic en \textbf{"Estudiantes"} o \textbf{"Usuarios"}
    \item Verás una tabla con todos los estudiantes de tu institución
    \columnas: Nombre, Cédula, Curso, Paralelo, Fecha test, Área predominante
    \item Puedes \textbf{filtar} por:
    \begin{itemize}
        \item Curso
        \item Paralelo
        \item Estado (activo/inactivo)
        \item Con/sin test completado
    \end{itemize}
\end{enumerate}

\subsection{Buscar un estudiante específico}
\begin{enumerate}
    \item Usa el \textbf{campo de búsqueda} (esquina superior derecha de la tabla)
    \item Escribe el nombre, cédula o apellido
    \item Los resultados se filtran automáticamente
\end{enumerate}

\subsection{Ver detalle de un estudiante}
\begin{enumerate}
    \item En la lista, haz clic en el \textbf{icono de ojo} (ver) junto al estudiante
    \item Se abrirá un modal con:
    \begin{itemize}
        \item Información personal completa
        \item Resultados del test (si lo completó)
        \item Gráfico personal
        \item Historial de tests (si hubo más de uno)
    \end{itemize}
\end{enumerate}

\section{Proceso 2: Generar Reporte Grupal Institucional}

\subsection{Reporte por curso/paralelo}
\begin{enumerate}
    \item Ve a \textbf{Reportes} $\rightarrow$ \textbf{Grupal}
    \item Selecciona los \textbf{filtros}:
    \begin{itemize}
        \item Institución: Tu institución (pre-seleccionada)
        \item Curso: Selecciona uno o múltiples
        \item Paralelo: Si aplica
        \item Fecha desde/hasta: Rango temporal
    \end{itemize}
    \item Elige el \textbf{formato de salida}:
    \begin{itemize}
        \item \textbf{PDF}: Para imprimir o presentar
        \item \textbf{Excel}: Para análisis detallado en hojas de cálculo
    \end{itemize}
    \item Haz clic en \textbf{"Generar Reporte"}
    \item El archivo se descargará automáticamente
\end{enumerate}

\textbf{Contenido del reporte grupal:}
\begin{itemize}
    \item Resumen ejecutivo con promedios
    \item Tabla detallada por estudiante
    \item Gráficos de distribución por área RIASEC
    \item Estadísticas por curso/paralelo
\end{itemize}

\section{Proceso 3: Seguimiento Individual}

\subsection{Ver reporte individual de un estudiante}
\begin{enumerate}
    \item En la lista de estudiantes, localiza al alumno
    \item Haz clic en el \textbf{botoón "Ver Reporte"} (icono de PDF)
    \item Se abrirá el PDF en una nueva pestaña
    \item Puedes:
    \begin{itemize}
        \item \textbf{Imprimir} para archivo físico
        \item \textbf{Descargar} a tu computadora
        \item \textbf{Enviar por email} (si está habilitado)
    \end{itemize}
\end{enumerate}

\textbf{El reporte individual incluye:}
\begin{itemize}
    \item Datos completos del estudiante
    \item Fecha del test
    \item Gráfico de resultados (radar y barras)
    \item Áreas predominantes (top 3)
    \item Lista de carreras recomendadas
    \item Código QR para validación oficial
\end{itemize}

\subsection{Realizar seguimiento}
\begin{enumerate}
    \item En el detalle del estudiante, haz clic en \textbf{"Agregar Nota"} o \textbf{Seguimiento"
    \item Escribe observaciones sobre:
    \begin{itemize}
        \item Progreso académico
        \item Intereses vocacionales declarados
        \item Acciones de orientación realizadas
        \item Recomendaciones adicionales
    \end{itemize}
    \item Guarda la nota
    \item Puedes ver el historial completo de notas
\end{enumerate}

\section{Exportar Datos a Excel}
\begin{enumerate}
    \item Ve a \textbf{Reportes} $\rightarrow$ \textbf{Exportar}
    \item Configura filtros (institución, curso, fecha)
    \item Selecciona \textbf{Excel (.xlsx)}
    \item Haz clic en \textbf{Exportar"}
    \item El archivo contiene:
    \begin{itemize}
        \item Datos crudos de todos los tests
        \item Columnas: Estudiante, Cédula, Curso, Paralelo, R,I,A,S,E,C scores, Fecha
        \item Formato listo para análisis en Excel
    \end{itemize}
\end{enumerate}

% ===========================================================================
% CAPÍTULO 5: GUÍA PARA ZONAL
% ===========================================================================
\chapter{GUÍA PARA ZONAL (\rolezonal)}

\section{¿Qué es un Zonal?}
Un \textbf{Coordinador Zonal} supervisa múltiples instituciones educativas dentro de una zona geográfica definida (ej: Zona Costa, Sierra, Amazonía, Insular).

\section{¿Qué puedes hacer como Zonal?}
\begin{itemize}
    \item Ver \textbf{todas las instituciones} de tu zona
    \item Generar reportes zonales agregados
    \item Filtrar por institución específica
    \item Comparar resultados entre instituciones
    \item Exportar datos masivos a Excel
    \item Descargar reportes individuales de cualquier estudiante de tu zona
\end{itemize}

\textbf{Limitación}: Solo puedes ver datos de instituciones asignadas a \underline{tu zona}. No puedes ver otras zonas.

\section{Acceso al Panel Zonal}
\begin{enumerate}
    \item Inicia sesión con credenciales de Zonal
    \item Serás redirigido al \textbf{Dashboard Zonal}
    \item O navega a \texttt{/admin/zona}
\end{enumerate}

\section{Dashboard Zonal - Vista General}
El dashboard zonal muestra:
\begin{itemize}
    \item \textbf{Mapa de zona}: Instituciones por ubicación (si está habilitado)
    \item \textbf{Estadísticas agregadas}:
    \begin{itemize}
        \item Total estudiantes en la zona
        \item Tests realizados
        \item Promedio general de puntajes RIASEC
    \end{itemize}
    \item \textbf{Gráfico comparativo}: Rendimiento por institución
    \item \textbf{Top áreas vocacionales} de la zona
    \item \textbf{Lista de instituciones} con enlaces rápidos
\end{itemize}

\section{Proceso 1: Reporte Zonal General}
\begin{enumerate}
    \item Ve a \textbf{Reportes} $\rightarrow$ \textbf{Zonal}
    \item Selecciona:
    \begin{itemize}
        \item Zona (tu zona por defecto)
        \item Rango de fechas (opcional)
    \end{itemize}
\item Elige formato: \textbf{PDF} o \textbf{Excel}
\item Haz clic en \textbf{Generar}
\item Descarga el archivo
\end{enumerate}

\textbf{Contenido del reporte zonal:}
\begin{itemize}
    \item Resumen de todas las instituciones de la zona
    \item Estadísticas consolidadas por categoría RIASEC
    \item Tabla comparativa entre instituciones
    \item Identificación de tendencias vocacionales regionales
    \item Gráficos de distribución
\end{itemize}

\section{Proceso 2: Reporte por Institución Específica}
\begin{enumerate}
    \item En el Dashboard, busca la \textbf{lista de instituciones}
    \item Haz clic en el \textbf{nombre} de la institución
    \item O usa el filtro \textbf{Institución} en Reportes
    \item Selecciona curso/paralelo si deseas filtrar más
    \item Genera reporte (PDF/Excel)
\end{enumerate}

\section{Comparar Instituciones}
\begin{enumerate}
    \item Ve a \textbf{Reportes} $\rightarrow$ \textbf{Comparativo}
    \item Selecciona \textbf{múltiples instituciones} (checkboxes)
    \item Elige período de tiempo
\item Haz clic en \textbf{Comparar}
\item Verás un informe con:
    \begin{itemize}
        \item Tabla comparativa de promedios por área
        \item Gráfico de barras agrupadas
        \item Diferencias significativas
    \end{itemize}
\end{enumerate}

% ===========================================================================
% CAPÍTULO 6: GUÍA PARA ADMINISTRADOR
% ===========================================================================
\chapter{GUÍA PARA ADMINISTRADOR (\roleadmin)}

\section{¿Qué puede hacer un Administrador?}
El Administrador tiene control total sobre el sistema:
\begin{itemize}
    \item \textbf{Gestión de usuarios}: Crear, editar, eliminar, asignar roles
    \item \textbf{Gestión de instituciones}: Alta, baja, modificación
    \item \textbf{Gestión de preguntas}: Banco de preguntas, importación masiva
    \item \textbf{Reportes globales}: Acceso a todos los datos del sistema
    \item \textbf{Configuración}: Parámetros del test, umbrales
    \item \textbf{Supervisión}: Monitoreo de actividad del sistema
\end{itemize}

\section{Acceso al Panel Admin}
\begin{enumerate}
    \item Inicia sesión con credenciales de administrador
    \item Serás redirigido al \textbf{Dashboard Administrativo}
    \item URL directa: \texttt{/admin} o \texttt{/admin/dashboard}
\end{enumerate}

\section{Dashboard Admin - Elementos}
\begin{table}[htbp]
\centering
\begin{tabularx}{\textwidth}{|l|X|}
\hline
\textbf{Widget} & \textbf{Información que muestra} \\
\hline
Estudiantes totales & Cantidad de usuarios con rol estudiante \\
\hline
Tests realizados & Número de tests completados hoy/esta semana/mes \\
\hline
Instituciones activas & Número de instituciones registradas \\
\hline
Distribución RIASEC & Gráfico de promedios globales por categoría \\
\hline
Actividad reciente & Últimos registros, tests, logs importantes \\
\hline
\end{tabularx}
\end{table}

\section{Proceso 1: Gestionar Instituciones}

\subsection{Crear nueva institución}
\begin{enumerate}
    \item Ve a \textbf{ Instituciones} $\rightarrow$ \textbf{Nueva Institución"}
    \item Completa el formulario:
    \begin{itemize}
        \item Nombre de la institución
        \item Código AMIE (código oficial del Ministerio de Educación)
        \item Tipo (Pública/Privada/FIS)
        \item Dirección completa
        \item Teléfono
        \item Email de contacto
        \item Representante legal
        \item Zona (Costal, Sierra, Amazónica, Insular)
    \end{itemize}
    \item Haz clic en \textbf{Guardar"}
    \item La institución aparece en la lista con estado \textbf{Activo"}
\end{enumerate}

\textbf{Importante}: El \textbf{código AMIE} es único y obligatorio. Verifica que no exista antes de guardar.

\subsection{Editar institución}
\begin{enumerate}
    \item En la lista de instituciones, haz clic en \textbf{"Editar"} (icono lápiz)
    \item Modifica los campos necesarios
    \item Haz clic en \textbf{Actualizar"}
\end{enumerate}

\subsection{Desactivar institución}
\begin{enumerate}
    \item En la lista, haz clic en \textbf{Desactivar"} (icono de interruptor)
    \item Confirma la acción
    \item La institución se marca como inactiva, pero \underline{no se eliminan sus datos}
\end{enumerate}

\section{Proceso 2: Gestionar Usuarios}

\subsection{Crear usuario (DecE, Zonal, Estudiante)}
\begin{enumerate}
    \item Ve a \textbf{Usuarios} $\rightarrow$ \textbf{Nuevo Usuario"}
    \item Completa el formulario:
    \begin{itemize}
        \item \textbf{Tipo de registro}: Manual o Masivo (Excel)
        \item \textbf{Rol}: Selecciona (estudiante, DECE, zonal, admin)
        \item \textbf{Cédula}: Número único
        \item \textbf{Nombre completo}
        \item \textbf{Email} (para login y notificaciones)
        \item \textbf{Contraseña temporal} (el usuario debe cambiarla)
        \item \textbf{Institución}: (solo para estudiantes y DECE)
        \item \textbf{Zona}: (solo para zonales)
    \end{itemize}
    \item Haz clic en \textbf{Crear Usuario"}
    \item \textbf{Nota}: Si es estudiante, se envía email de bienvenida (si SMTP configurado)
\end{enumerate}

\subsection{Importar usuarios masivamente (Excel)}
\begin{enumerate}
    \item Ve a \textbf{Usuarios} $\rightarrow$ \textbf{Importar"}
    \item Descarga la \textbf{plantilla Excel} oficial
    \item Llena la plantilla con estos campos:
    \begin{itemize}
        \item cedula
        \item nombre
        \item email
        \item password (temporal)
        \item rol
        \item institucion\_id (si aplica)
        \item zona (si aplica)
    \end{itemize}
    \item Guarda el archivo
    \item En la página de importación, haz clic en \textbf{Seleccionar archivo"}
    \item Elige tu Excel y haz clic en \textbf{Importar"}
    \item El sistema mostrará un resumen: X usuarios creados, Y errores, Z duplicados
\end{enumerate}

\underline{Formato de plantilla Excel:}
\begin{lstlisting}[basicstyle=\ttfamily\small]
cedula | nombre | email | password | rol | institucion_id | zona
1234567| Juan Pérez| juan@mail.com| Temp123| estudiante| 5| sierra
\end{lstlisting}

\subsection{Editar usuario}
\begin{enumerate}
    \item En la lista de usuarios, busca el usuario
    \item Haz clic en \textbf{"Editar"} (icono lápiz)
    \item Puedes modificar:
    \begin{itemize}
        \item Datos personales (nombre, email, cédula - con cuidado)
        \item Rol (cambiar permiso)
        \item Institución asignada (para DECE/estudiantes)
        \item Zona (para zonales)
        \item Estado (activo/inactivo)
        \item Forzar cambio de contraseña (checkboxes)
    \end{itemize}
    \item Haz clic en \textbf{Guardar cambios"}
\end{enumerate}

\subsection{Eliminar usuario}
\begin{enumerate}
    \item En la lista, marca el checkbox del usuario
    \item Haz clic en \textbf{Eliminar seleccionados"}
    \item \textbf{¡PRECAUCIÓN!} Confirma que realmente deseas eliminarlo
    \item El usuario se elimina permanentemente (o se marca como inactivo, según configuración)
\end{enumerate}

\section{Proceso 3: Gestionar Banco de Preguntas}

\subsection{Ver preguntas existentes}
\begin{enumerate}
    \item Ve a \textbf{Preguntas} en el menú
    \item Verás una tabla con todas las preguntas
    \item Columnas: ID, Texto, Categoría (R,I,A,S,E,C), Tipo (intereses/habilidades/valores), Estado
    \item Puedes filtrar por:
    \begin{itemize}
        \item Categoría
        \item Tipo
        \item Estado (activas/inactivas)
    \end{itemize}
\end{enumerate}

\subsection{Agregar pregunta individual}
\begin{enumerate}
    \item Haz clic en \textbf{Nueva Pregunta"}
    \item Completa:
    \begin{itemize}
        \item \textbf{Texto}: La pregunta (clara, sin errores ortográficos)
        \item \textbf{Categoría RIASEC}: R, I, A, S, E, o C
        \item \textbf{Tipo}: Intereses, Habilidades o Valores
        \item \textbf{Orden}: Posición en que aparecerá (0 por defecto)
        \item \textbf{Estado}: Activo / Inactivo
    \end{itemize}
    \item Haz clic en \textbf{Guardar"}
\end{enumerate}

\subsection{Importar preguntas masivamente (Excel)}
\begin{enumerate}
    \item Haz clic en \textbf{Importar Preguntas"}
    \item Descarga la \textbf{plantilla oficial} (Excel .xlsx)
    \item La plantilla tiene estas columnas:
    \begin{lstlisting}[basicstyle=\ttfamily\small]
categoria | tipo | pregunta | peso
R         | intereses| Me gusta trabajar con herramientas| 1
    \end{lstlisting}
    \item Llena tantas filas como preguntas tengas
    \item Guarda el archivo
    \item En la página de importación, haz clic en \textbf{Seleccionar archivo}
    \item Arrastra tu archivo Excel
    \item Verás una \textbf{previsualización} de las primeras 5 preguntas
    \item Si se ve correcto, haz clic en \textbf{Importar"}
    \item El sistema procesará y mostrará: X importadas, Y duplicadas, Z con errores
\end{enumerate}

\textbf{Formatos soportados:}
\begin{itemize}
    \item \textbf{Excel (.xlsx, .xls)}: Recomendado
    \item \textbf{Word (.docx, .doc)}: Formato estructurado
\end{itemize}

\subsection{Eliminar preguntas}
\begin{enumerate}
    \item En la lista de preguntas, marca los checkboxes
    \item Haz clic en \textbf{Eliminar seleccionadas"}
    \item Confirma la eliminación
    \item \textbf{¡CUIDADO!}: Las preguntas eliminadas se quitan del test
\end{enumerate}

\subsection{Cambiar estado de preguntas}
\begin{enumerate}
    \item En la lista, usa el \textbf{selector de estado} (Activo/Inactivo)
    \item Las preguntas \textbf{inactivas} no aparecen en el test
    \item Útil para desactivar preguntas obsoletas sin borrarlas
\end{enumerate}

\section{Proceso 4: Generar Reportes Globales}

\subsection{Reporte individual de cualquier estudiante}
\begin{enumerate}
    \item Ve a \textbf{Reportes} $\rightarrow$ \textbf{Individual"}
    \item Busca el estudiante por:
    \begin{itemize}
        \item Cédula
        \item Nombre
        \item Institución + curso
    \end{itemize}
    \item Selecciona al estudiante
    \item Haz clic en \textbf{Generar PDF"}
    \item Descarga el reporte
\end{enumerate}

\subsection{Reporte grupal (todos los datos)}
\begin{enumerate}
    \item Ve a \textbf{Reportes} $\rightarrow$ \textbf{Grupal"}
    \item Aplica \textbf{filtros} (opcional):
    \begin{itemize}
        \item Zona
        \item Institución (una o múltiples)
        \item Curso
        \item Paralelo
        \item Rango de fechas
    \end{itemize}
    \item Selecciona \textbf{formato}: PDF o Excel
    \item Haz clic en \textbf{Generar"}
    \item Descarga el archivo
\end{enumerate}

\textbf{Reporte Excel contiene:}
\begin{itemize}
    \item Hoja 1: Datos de estudiantes
    \item Hoja 2: Promedios por categoría RIASEC
    \item Hoja 3: Estadísticas por institución
    \item Hoja 4: Top carreras recomendadas
\end{itemize}

\subsection{Estadísticas en tiempo real}
\begin{itemize}
    \item En el Dashboard, todos los gráficos son \textbf{en tiempo real}
    \item Se actualizan cada vez que recargas la página
    \item Reflejan los datos más recientes de la base de datos
\end{itemize}

\section{Proceso 5: Configuración del Sistema}

\subsection{Ajustar parámetros del test}
\begin{enumerate}
    \item Ve a \textbf{Configuración} $\rightarrow$ \textbf{Parámetros"}
    \item Puedes modificar:
    \begin{itemize}
        \item \textbf{Tiempo entre tests}: Meses que debe pasar para repetir (default: 6)
        \item \textbf{Umbral APTO}: Porcentaje mínimo (default: 80\%)
        \item \textbf{Umbral POTENCIAL}: Porcentaje mínimo (default: 60\%)
        \item \textbf{Preguntas por categoría}: Número de preguntas (default: 10)
    \end{itemize}
    \item \textbf{¡CUIDADO!}: Cambios aquí afectan a \underline{todos los usuarios}
    \item Haz clic en \textbf{Guardar configuración"}
\end{enumerate}

\subsection{Gestionar zonas geográficas}
\begin{enumerate}
    \item Ve a \textbf{Configuración} $\rightarrow$ \textbf{Zonas"}
    \item Puedes ver/cambiar:
    \begin{itemize}
        \item Zona Costal
        \item Zona Sierra
        \item Zona Amazónica
        \item Zona Insular
    \end{itemize}
    \item Asignar instituciones a zonas
\end{enumerate}

% ===========================================================================
% CAPÍTULO 7: USABILIDAD Y CONSEJOS
% ===========================================================================
\chapter{USABILIDAD Y CONSEJOS PRÁCTICOS}

\section{Consejos Generales para Todos los Usuarios}
\begin{tcolorbox}[colback=lightgray, colframe=successgreen, title=Mejores Prácticas]
\begin{itemize}
    \item \textbf{No uses el botón "Atrás" del navegador}. Usa los botones del sistema (« Anterior)
    \item \textbf{No cierres el navegador} a mitad del test sin guardar
    \item \textbf{Usa contraseñas seguras}: combinación de letras, números y símbolos
    \item \textbf{Cierra sesión} en computadoras compartidas
    \item \textbf{Reporta errores}: Si ves algo inusual, avisa al administrador
\end{itemize}
\end{tcolorbox}

\section{Navegación por el Sistema}
\begin{itemize}
    \item \textbf{Menú superior}: Acceso rápido a perfil y logout
    \item \textbf{Menú lateral}: Funcionalidades principales (varía por rol)
    \item \textbf{Breadcrumbs}: Muestra tu ubicación actual (ej: Inicio > Reportes > Individual)
    \item \textbf{Botones de acción}: Siempre verdes o azules (positivos)
    \item \textbf{Botones de cancelar}: Grises o rojos (negativos)
\end{itemize}

\section{Interpretación de Iconos}
\begin{table}[htbp]
\centering
\begin{tabularx}{\textwidth}{|c|l|}
\hline
\textbf{Icono} & \textbf{Significado} \\
\hline
\includegraphics[height=0.8em]{eye.png} & Ver detalle \\
\hline
\includegraphics[height=0.8em]{edit.png} & Editar \\
\hline
\includegraphics[height=0.8em]{delete.png} & Eliminar \\
\hline
\includegraphics[height=0.8em]{download.png} & Descargar \\
\hline
\includegraphics[height=0.8em]{print.png} & Imprimir \\
\hline
\includegraphics[height=0.8em]{pdf.png} & PDF \\
\hline
\includegraphics[height=0.8em]{excel.png} & Excel \\
\hline
\includegraphics[height=0.8em]{search.png} & Buscar \\
\hline
\end{tabularx}
\caption{Iconos comunes en la interfaz}
\end{table}

\section{Mensajes del Sistema}
El sistema muestra 4 tipos de mensajes:

\begin{tcolorbox}[colback=successgreen!10, colframe=successgreen, title=ÉXITO (Verde)]
Operación completada correctamente.\\
Ej: "Usuario creado exitosamente"
\end{tcolorbox}

\begin{tcolorbox}[colback=dangerred!10, colframe=dangerred, title=ERROR (Rojo)]
Algo salió mal. La acción no se completó.\\
Ej: "Contraseña incorrecta"
\end{tcolorbox}

\begin{tcolorbox}[colback=infoblue!10, colframe=infoblue, title=INFO (Azul)]
Información importante o sugerencias.\\
Ej: "Tu test expira en 3 días"
\end{tcolorbox}

\begin{tcolorbox}[colback=warningorange!10, colframe=warningorange, title=ADVERTENCIA (Amarillo)]
Precaución, algo no es ideal pero no es error.\\
Ej: "Estás por eliminar un registro permanentemente"
\end{tcolorbox}

% ===========================================================================
% CAPÍTULO 8: PREGUNTAS FRECUENTES (FAQ)
% ===========================================================================
\chapter{PREGUNTAS FRECUENTES (FAQ)}

\section{Estudiantes}

\subsection{No puedo iniciar sesión}
\textbf{P: Al intentar login, dice "Usuario o contraseña incorrectos"}\\
\textbf{R:}
\begin{enumerate}
    \item Verifica que estás usando el \textbf{usuario correcto} (cédula o email)
    \item Asegúrate de que \underline{Caps Lock no esté activado}
    \item Si olvidaste tu contraseña, haz clic en \textbf{¿Olvidaste tu contraseña?"}
    \item Si el problema persiste, contacta a tu DECE o administrador
\end{enumerate}

\subsection{El test se cerró accidentalmente}
\textbf{P: Cerré la ventana a mitad del test, ¿puedo recuperar mis respuestas?}\\
\textbf{R:} Sí:
\begin{enumerate}
    \item Vuelve a la página del test
    \item Haz clic en \textbf{Recuperar respuestas guardadas"}
    \item Verás las preguntas que ya respondiste
    \item Continúa desde donde te quedaste
\end{enumerate}
\textbf{Importante}: Esto solo funciona si no enviaste el test final. Una vez finalizado, no se puede modificar.

\subsection{Veo resultados antiguos}
\textbf{P: Realicé un test hace 7 meses, ¿puedo hacer uno nuevo?}\\
\textbf{R:} Sí. Deben haber pasado al menos \underline{6 meses (180 días)}. Si han pasado 7, ya puedes realizar un nuevo test. El sistema calculará automáticamente el nuevo resultado y guardará ambos (historial).

\subsection{No entiendo mis resultados}
\textbf{P: ¿Qué significa que tengo "APTO" en Realista?}\\
\textbf{R:}
\begin{itemize}
    \item \textbf{APTO} (verde): Tu afinidad con esa área es alta (≥80\%). Es una fortaleza.
    \item \textbf{POTENCIAL} (amarillo): Tienes buen nivel (60-79\%) pero puedes fortalecerlo.
    \item \textbf{POR REFORZAR} (rojo): Es un área de menor afinidad (<60\%).
\end{itemize}
El \textbf{área predominante} es la que tiene puntaje más alto. Enfócate en carreras relacionadas con esa área.

\section{DECE y Zonal}

\subsection{No veo estudiantes de cierta institución}
\textbf{P: En mi lista faltan estudiantes de la escuela "X"}\\
\textbf{R:}
\begin{enumerate}
    \item Verifica que la institución esté \textbf{asignada a tu cuenta}
    \item Como DECE, solo ves tu institución
    \item Como Zonal, debes tener la zona correcta configurada
    \item Contacta al Administrador para que asigne la institución a tu rol
\end{enumerate}

\subsection{Reporte grupal no muestra todos los datos}
\textbf{P: Generé un reporte y solo aparecen 10 de 30 estudiantes}\\
\textbf{R:}
\begin{itemize}
    \item Revisa los \textbf{filtros aplicados} (curso, paralelo, fecha)
    \item Asegúrate de que el rango de fechas incluya el período deseado
    \item Verifica que los estudiantes hayan completo el test (solo aparecen con resultados)
    \item Intenta con \textbf{"Limpiar filtros"} y generar de nuevo
\end{itemize}

\subsection{No puedo importar preguntas}
\textbf{P: Al importar Excel, dice "Formato inválido"}\\
\textbf{R:}
\begin{enumerate}
    \item Verifica que uses la \textbf{plantilla oficial} (no modifiques encabezados)
    \item Las columnas deben ser exactamente: categoria, tipo, pregunta, peso
    \item La categoría debe ser una letra: R, I, A, S, E, o C
    \item El tipo: intereses, habilidades o valores
    \item El archivo debe ser .xlsx (no .xls en algunas versiones)
    \item Tamaño máximo: 10MB
\end{enumerate}

\section{Administradores}

\subsection{Error al crear usuario}
\textbf{P: Al crear usuario manual, dice "Email ya existe"}\\
\textbf{R:} El email debe ser único. Verifica:
\begin{enumerate}
    \item Que el email no esté ya registrado (busca en la lista)
    \item Si es un estudiante que ya tenía cuenta, actívala en lugar de crearla nueva
    \item Usa un email diferente (puedes agregar +númerosin Gmail)
\end{enumerate}

\subsection{Reporte PDF no se genera}
\textbf{P: Hago clic en "Generar PDF" y no descarga nada}\\
\textbf{R:}
\begin{itemize}
    \item Verifica que \textbf{TCPDF esté instalado} (composer install)
    \item Revisa logs de PHP: \texttt{/var/log/php\_errors.log}
    \item Asegúrate de que el estudiante tenga resultados guardados
    \item Prueba con otro navegador
    \item Si persiste, reinicia el servidor Apache/PHP
\end{itemize}

\section{Todos los Usuarios}

\subsection{La página se queda cargando}
\textbf{P: Hago clic en algo y aparece la rueda girando eternamente}\\
\textbf{R:}
\begin{enumerate}
    \item Espera 30 segundos (a veces el servidor es lento)
    \item Recarga la página (F5 o Ctrl+R)
    \item Verifica tu conexión a internet
    \item Si es un test, revisa si guardó respuestas (localStorage)
    \item Contacta al administrador si es recurrente
\end{enumerate}

\subsection{No veo gráficos/estadísticas}
\textbf{P: Las páginas de resultados están en blanco o sin gráficos}\\
\textbf{R:}
\begin{itemize}
    \item Asegúrate de que JavaScript esté \textbf{habilitado} en tu navegador
    \item Prueba con Chrome o Firefox
    \item Limpia la caché del navegador (Ctrl+Shift+Del)
    \item Verifica que Chart.js esté cargando (F12 $\rightarrow$ Console)
\end{itemize}

% ===========================================================================
% CAPÍTULO 9: GLOSARIO
% ===========================================================================
\chapter{GLOSARIO DE TÉRMINOS}

\begin{description}[leftmargin=2cm, font=\normalfont]
    \item[RIASEC] Modelo tipológico de John Holland: Realista, Investigador, Artístico, Social, Emprendedor, Convencional.
    
    \item[Escala Likert] Escala de medición ordinal del 1 al 5 que mide grado de acuerdo/desacuerdo.
    
    \item[DECE] Departamento de Estudios y Consejería Estudiantil. Personal orientador en instituciones educativas.
    
    \item[Zonal] Coordinador zonal que supervisa múltiples instituciones en una región geográfica.
    
    \item[Código AMIE] Código único que el Ministerio de Educación asigna a cada institución educativa.
    
    \item[Paralelo] División o grupo de estudiantes dentro de un mismo curso (ej: 1ro "A", 1ro "B").
    
    \item[PDF] Portable Document Format. Formato de archivo para documentos.
    
    \item[Excel] Hoja de cálculo de Microsoft (formato .xlsx).
    
    \item[LocalStorage] Almacenamiento del navegador que guarda datos temporalmente.
    
    \item[QR] Código QR (Quick Response) que contiene información escaneable.
    
    \item[Middleware] Software que filtra peticiones antes de llegar al controlador (autenticación).
    
    \item[API] Application Programming Interface (no expuesta directamente a usuarios).
    
    \item[SMTP] Protocolo para envío de correos electrónicos.
    
    \item[Roles] Niveles de acceso: Administrador > Zonal > DECE > Estudiante.
    
    \item[Test de Re-take] Segundo test vocacional (permitido cada 6 meses).
    
    \item[Perfil vocacional] Conjunto de puntajes RIASEC que definen tendencias de carrera.
\end{description}

% ===========================================================================
% CAPÍTULO 10: SOPORTE Y CONTACTO
% ===========================================================================
\chapter{SOPORTE Y CONTACTO}

\section{¿Necesitas ayuda?}
Si tienes problemas o preguntas:

\begin{itemize}
    \item \textbf{Primero}: Consulta este manual (índice al inicio)
    \item \textbf{Segundo}: Revisa la sección \textbf{Preguntas Frecuentes} (Capítulo 8)
    \item \textbf{Tercero}: Contacta al administrador de tu institución (DECE)
    \item \textbf{Cuarto}: Si es un error técnico, envía un ticket a \textbf{soporte@testvocacional.edu.ec}
\end{itemize}

\section{Información para Reportar Errores}
Al reportar un error, incluye:
\begin{enumerate}
    \item \textbf{Tu rol}: Estudiante / DECE / Zonal / Administrador
    \item \textbf{Institución}: Nombre de tu escuela/colegio
    \item \textbf{Paso a paso}: Qué hiciste antes del error
    \item \textbf{Descripción exacta}: Copia y pega mensajes de error
    \item \textbf{Navegador}: Chrome/Firefox y versión
    \item \textbf{Fecha y hora} del incidente
    \item \textbf{Captura de pantalla} (si es posible)
\end{enumerate}

\section{Contacto Institucional}
\begin{table}[htbp]
\centering
\begin{tabularx}{\textwidth}{|l|X|}
\hline
\textbf{Concepto} & \textbf{Contacto} \\
\hline
Soporte técnico general & soporte@testvocacional.edu.ec \\
Coordinación DECE & dece@testvocacional.edu.ec \\
Administrador del sistema & admin@testvocacional.edu.ec \\
\hline
\end{tabularx}
\end{table}

\section{Tiempos de Respuesta}
\begin{itemize}
    \item \textbf{Errores críticos}: 2-4 horas (sistema caído, sin acceso)
    \item \textbf{Problemas funcionales}: 8-24 horas (no puede realizar test)
    \item \textbf{Consultas generales}: 1-2 días hábiles
\end{itemize}

\section{Cursos de Capacitación}
Se ofrecen capacitaciones virtuales periódicas:
\begin{itemize}
    \item \textbf{Para DECE}: "Gestión del Test Vocacional en tu Institución"
    \item \textbf{Para Zonales}: "Reportes y Análisis de Datos Vocacionales"
    \item \textbf{Para Administradores}: "Administración Avanzada del Sistema"
\end{itemize}
Contacta a \textbf{capacitacion@testvocacional.edu.ec} para inscribirte.

% ===========================================================================
% APÉNDICES
% ===========================================================================
\appendix
\chapter{APÉNDICES}

\section{Checklist - Antes del Test (Estudiante)}
\begin{itemize}
    \item Tener acceso a una computadora (no celular)
    \item Conexión a internet estable
    \item Navegador actualizado (Chrome recomendado)
    \item 45-60 minutos de tiempo disponible (sin interrupciones)
    \item Ambientes tranquilo y bien iluminado
    \item Lápiz y papel para notas (opcional)
    \item No consumir cafeína o azúcar en exceso antes del test
\end{itemize}

\section{Checklist - DECE: Revisión de Resultados}
\begin{itemize}
    \item Verificar que todos los estudiantes de la institución hayan sido asignados
    \item Revisar estudiantes pendientes (que no han hecho test)
    \item Analizar gráficos de distribución vocacional
    \item Identificar estudiantes con APTO en áreas clave
    \item Generar reporte grupal por curso
    \item Planificar sesiones de orientación individual
    \item Documentar observaciones en el sistema
\end{itemize}

\section{Checklist - Administrador: Mantenimiento Mensual}
\begin{itemize}
    \item Revisar logs de errores
    \item Verificar backups automáticos de BD
    \item Actualizar preguntas (agregar nuevo, eliminar obsoletas)
    \item Revisar usuarios inactivos (limpiar si es necesario)
    \item Monitorear espacio en disco
    \item Verificar envío de emails (SMTP)
    \item Actualizar lista de instituciones
\end{itemize}

\section{Atajos de Teclado}
\begin{table}[htbp]
\centering
\begin{tabularx}{\textwidth}{|l|l|}
\hline
\textbf{Atajo} & \textbf{Acción} \\
\hline
F5 & Recargar página \\
\hline
Ctrl+F & Buscar en la página \\
\hline
Ctrl+P & Imprimir página actual \\
\hline
Ctrl+S & Guardar (en formularios) \\
\hline
Alt + Flecha Izq & Volver atrás (historial navegador) \\
\hline
\end{tabularx}
\end{table}

\section{Resolución de Problemas Rápidos}
\begin{table}[htbp]
\centering
\begin{tabularx}{\textwidth}{|l|X|}
\hline
\textbf{Síntoma} & \textbf{Solución inmediata} \\
\hline
No puedo loguearme & 1. Limpiar caché del navegador \\
& 2. Probar en modo incógnito \\
& 3. Recuperar contraseña \\
\hline
Test se cuelga & 1. Esperar 30 segundos \\
& 2. Recuperar respuestas guardadas \\
& 3. Si no funciona, reportar \\
\hline
PDF no descarga & 1. Verificar pop-ups habilitados \\
& 2. Probar navegador diferente \\
& 3. Desactivar ad-blocker \\
\hline
Gráficos no se ven & 1. Habilitar JavaScript \\
& 2. Actualizar navegador \\
& 3. Limpiar caché \\
\hline
\end{tabularx}
\end{table}

\section{Plantillas y Formatos}
Adjuntos a este manual (archivos separados):
\begin{itemize}
    \item \textbf{Plantilla Excel para importar preguntas} (preguntas\_template.xlsx)
    \item \textbf{Plantilla Excel para importar usuarios} (usuarios\_template.xlsx)
    \item \textbf{Formato de reporte grupal} (ejemplo\_reporte.pdf)
    \item \textbf{Glosario RIASEC} (carreras\_por_categoria.pdf)
\end{itemize}

\end{document}
